\documentclass[10pt,a4paper]{amsart}
\usepackage[margin=19mm]{geometry}
\usepackage{amsmath,amssymb,mathtools}
\usepackage[scr=rsfs,cal=euler]{mathalfa}
\usepackage{xfrac}
\usepackage[unicode,hidelinks,bookmarks=false]{hyperref}
\newcommand{\ie}{i.e.\ }
\newcommand{\cf}{cf.\ }
\newcommand{\resp}{respectively\ }
\newcommand{\Subsection}{\subsection}
\providecommand{\nobreakdash}{\nobreak\hbox{-}}
\setlength{\emergencystretch}{2em}
\multlinegap0pt
\usepackage{bm}
\usepackage{bbm}
\usepackage{dsfont}
\usepackage{xspace}
\usepackage{cancel}
\usepackage{mathtools}
\usepackage{amsthm}
\makeatletter
\@ifundefined{Thm}{\newtheorem{Thm}{Thm}}{}
\@ifundefined{Prop}{\newtheorem{Prop}[Thm]{Proposition}}{}
\makeatother
\makeatletter
\expandafter\ifx\csname psmatrix\endcsname\relax
\newenvironment{psmatrix}
  {\left(\begin{smallmatrix}}
\fi
\newcommand{\zz}{\mathbb Z}
\makeatother
\title[Weyl algebras on Braverman-Kazhdan spaces]{Weyl algebras on Braverman-Kazhdan spaces}
\author{Chun-Hsien Hsu}
\date{Source: arXiv:2606.01206v1 (CC BY 4.0)}
\begin{document}
\maketitle

\noindent\emph{Source and licence.} Weyl algebras on Braverman-Kazhdan spaces. Version arXiv:2606.01206v1 (CC BY 4.0), distributed under the Creative Commons Attribution 4.0 International licence, as recorded in the arXiv licence metadata for this exact version and shown on the abstract page https://arxiv.org/abs/2606.01206v1. Full rights notice: licenses/arxiv-2606.01206-CC-BY-4.0.txt.\par\smallskip

\noindent\textit{Editorial scope.} Counted unit: the Prop whose source label is \texttt{None} in the article (source0.tex lines 401--407, source label \texttt{None}), delivered together with the article's own complete original proof of it (source0.tex lines 409--440, one \texttt{proof} environment of 32 lines). No mathematics is written, repaired or re-derived: statement and proof are the article's own text line for line. Required context retained uncounted: none. Every definition, display and label of those lines is retained unchanged. Omitted from this package and not redistributed: 1--400, 408--408, 441--2076. Nothing inside a retained excerpt is omitted or altered: every retained body is delivered exactly as the article's lines are written, so no omission receipt of any kind accompanies this package. Citations: the retained excerpts carry no citation command at all, so no bibliography entry of the article is transcribed and no reference is redistributed. Reference adaptation: none was needed and none was made; every reference inside the delivered unit resolves inside it, so no unresolved reference or citation is produced. Printed slips kept literal: no printed word, symbol, hypothesis or number of the counted statement or of its proof was altered. Layout and class: the article's own class is \texttt{amsart} with packages including amscd, amsfonts, amssymb, mathtools, stmaryrd, geometry, fancybox, pdflscape, rotating, xy, enumerate, mathrsfs, color, hyperref; the wrapper is the lane's pinned \texttt{amsart} editorial wrapper at 10pt on A4, which restates the theorem-like environments the retained text needs and adds the article's own macro definitions that the retained text needs (\texttt{\textbackslash zz}), with extra wrapper packages (bm, bbm, dsfont, xspace, cancel, mathtools). The delivered numbering is the unit's own. Assets: no retained span contains a figure environment, a graphics-inclusion command, a PDF-page inclusion, a table or a TikZ picture; no image file is copied into this package, the package declares no asset dependency and no image file is redistributed with it. Licence: version arXiv:2606.01206v1 of this article is distributed under the Creative Commons Attribution 4.0 International licence, as recorded in the arXiv licence metadata for this exact version and shown on the abstract page https://arxiv.org/abs/2606.01206v1. A full rights notice is delivered at licenses/arxiv-2606.01206-CC-BY-4.0.txt. Characters: every retained excerpt is pure ASCII, so the delivered engine, XeLaTeX with its default Latin Modern text font, reports no missing character.
\begin{Prop}
\begin{enumerate}
    \item $W_{X_P}$ is a domain.
    \item The group of units in $W_{X_P}$ is $K^\times$.
    \item The center of $W_{X_P}$ is the field $K$.
\end{enumerate}
\end{Prop}

\begin{proof}
For $D\in W_{X_P}$ nonzero, choose a representative
\begin{align*}
    D=\sum_{\lambda\in I_D} p_{\lambda,D}(\mathbf{x}) \frac{\partial^{|\lambda|}}{\partial \mathbf{x}^\lambda}
\end{align*}
in the classical Weyl algebra $W_{V_P}$, where $I_D\subset \zz_{\ge 0}^{\dim V_P}$ is a (nonempty) finite subset and $p_{\lambda,D }\neq 0$ in $K[X_P]$ for all $\lambda\in I_D$.

For (1), as $W_{X_P}$ is a left primitive ring and hence a prime ring, it suffices to show $W_{X_P}$ is reduced. Suppose on the contrary that $D^2=0$ in $W_{X_P}$. Let
\begin{align*}
    \sum_{\lambda\in I} p_{\lambda,D^2}(\mathbf{x})\frac{\partial^{|\lambda|}}{\partial \mathbf{x}^\lambda}
\end{align*}
be the square of the representative of $D$ in $W_{V_P}$. Then $p_{\lambda,D^2}=0$ in $K[X_P]$. On the other hand, there is $\lambda\in I_D$ such that $|\lambda|$ is maximal, $2\lambda\in I$, and $2\lambda\neq \lambda_1+\lambda_2$ for any $\lambda_1,\lambda_2\in I_D$ unless $\lambda_1=\lambda_2=\lambda$.
%One can choose $\lambda$ successively as follows: choose $\lambda\in I_D$ which has an coordinate that is smallest among all other cooridnates of every $\lambda'$s. Fix this coordinate, which we may assume to be $1$. Let $I'$ be the set that consist of all $\lambda\in I_D$ collecting all $\lambda$ above. Then consider all $\lambda\in I'$ which has the smallest two coordinates... and keep going on. The process terminates when the set we consider is a singleton.
Then, $p_{2\lambda, D^2}=p_{\lambda,D}^2\neq 0$ in $K[X_P]$, which is a contradiction.

%Choose a $\lambda\in I_{D^2}$ with %$|\lambda|$ minimal. Then 
%\begin{align*}     0=D^2(\mathbf{x}^{\lambda})=\left(\prod_{i=1}^{\dim V_P} \lambda_i! \right)p_{\lambda,D^2}
%\end{align*}
%in $K[X_P]$. Thus $p_{\lambda,D^2}=0$ and %by an inductive argument it follows that %all polynomial coefficients are zero in %$K[X_P]$.



For (2), suppose $D\not\in K^\times$ and there is $D'\in W_{X_P}$ such that $DD'=1$. Then $D'\not\in K^\times.$ Let
\begin{align*}
    \sum_{\lambda\in I} p_{\lambda,DD'}(\mathbf{x})\frac{\partial^{|\lambda|}}{\partial \mathbf{x}^\lambda}
\end{align*}
be the multiplication of the representatives of $D$ and $D'$ in $W_{V_P}$, so $p_{0,DD'}=1$ and  $p_{\lambda,DD'}$ are zero in $K[X_P]$ for $\lambda\neq 0$. However, there are $\lambda\in I_D,\lambda'\in I_{D'}$ both nonzero such that $\lambda+\lambda'\in I$ and   $p_{\lambda+\lambda',DD'}=p_{\lambda,D}p_{\lambda',D'}\neq 0$
is nonzero in $K[X_P]$. Therefore, such $D'$ does not exist. 

Finally for (3) suppose $D$ lies in the center of $W_{X_P}$. Let $\lambda\in I_D$ with $|\lambda|$ minimal. Assume first $|\lambda|>0$. Choose $i$ such that $\lambda_i\neq 0$. Let $e_i\in \zz_{\ge 0}^{\mathrm{dim}V_P}$ be the vector whose $i$th entry is $1$ and zero otherwise. Then $p_{\lambda-e_i,[D,x_i]}=\lambda_i p_{\lambda,D}\neq 0$
in $K[X_P]$, contradicting the assumption that $[D,x_i]=0$. Suppose $\lambda=0$ and $p_0$ is nonconstant. Then there is $i$ such that $\partial_i(p_0)\neq 0$ in $K[X_P]$, and thus $[\partial_i,p_0]\neq 0$, a contradiction.
\end{proof}

\end{document}
